\providecommand{\key}[1]{\mbox{\fbox{\rule{0pt}{1.9ex}\hspace{0.15em}#1\hspace{0.15em}}}\hspace{0.1em}}
\providecommand{\kpl}{\key{$+$}}\providecommand{\kmi}{\key{$-$}}\providecommand{\kpo}{\key{$\times$}}\providecommand{\kdi}{\key{$\div$}}\providecommand{\kis}{\key{$=$}}

\paragrafos{Υπολογιστής τσέπης: ένα εργαλείο, όχι ένα μυαλό}{Α.3.3, «Υπολογισμοί με τη βοήθεια υπολογιστή τσέπης»}

\begin{anagke}
Μετά το παζάρι, το ταμείο έχει γεμίσει αποδείξεις. Ο Νίκος, που κρατάει το ταμείο, αρχίζει να τις προσθέτει με το χέρι και βγάζει $68{,}79$ €. Η Πολυξένη τις ξαναπροσθέτει με το κομπιουτεράκι του κινητού της και βγάζει $77{,}79$ €. «Ένας από τους δύο έχει λάθος», λέει ο Ηλίας. «Το κομπιουτεράκι δεν κάνει λάθος», απαντά ο Βασίλης. «Άρα ο Νίκος».

Η Πολυξένη κοιτάζει την οθόνη της και βλέπει ότι πάτησε κατά λάθος το $23{,}45$ ως $32{,}45$: τα ψηφία αντιστράφηκαν, και η διαφορά στο σύνολο είναι ακριβώς $9$ ευρώ. «Το κομπιουτεράκι έκανε σωστή πράξη με λάθος δεδομένα», παρατηρεί η κυρία Αλεξάνδρα. «Δεν σκέφτεται: κάνει ό,τι του πεις». Και ο Στέλιος, που θέλει να δοκιμάσει το εργαλείο, πατάει $1:3$ και βλέπει $0{,}33333333$ με οκτώ ψηφία. «Χάλασε; Το $\frac{1}{3}$ δεν είναι ακριβώς αυτό!».

Στο κεφάλαιο αυτό η παρέα μαθαίνει να χρησιμοποιεί το κομπιουτεράκι σωστά: πώς πατάμε τις πράξεις, πώς φυλάμε ενδιάμεσα αποτελέσματα στη μνήμη, πώς ξέρουμε αν η οθόνη λέει αλήθεια.
\end{anagke}

\begin{lexi}{υπολογιστής, μνήμη, οθόνη}{υπολογιστής, μνήμη}
Ο \lexiword{υπολογιστής} είναι αυτός που «υπολογίζει», από το «λόγος, λογαριασμός»: το να υπολογίζεις σημαίνει να κάνεις λογαριασμό. Παλιά οι υπολογιστές ήταν άνθρωποι με άβακα. Ο \lexiword{άβακας} ήταν ένας πίνακας με χάντρες. Σήμερα το μηχάνημα έχει δική του \lexiword{μνήμη}, ένα «συρτάρι» όπου φυλάμε έναν αριθμό για να τον ξαναβρούμε αργότερα.
\end{lexi}

\begin{metafrasi}
Πρόσθεσε $12{,}35$ και $7{,}8$. & \key{12,35} \kpl\ \key{7,8} \kis\\
Φύλαξε τον αριθμό της οθόνης στη μνήμη. & \key{M+}\\
Αφαίρεσε από τη μνήμη τον αριθμό της οθόνης. & \key{M$-$}\\
Δείξε μου τι υπάρχει στη μνήμη. & \key{MR}\\
Άδειασε τη μνήμη. & \key{MC}\\
Πρώτα η παρένθεση, μετά το υπόλοιπο. & \key{(}\ldots\key{)} ή ενδιάμεσα αποτελέσματα\\
\end{metafrasi}

\tmima{Η θεωρία}

\tmimasmall{Οι τέσσερις πράξεις}

Όλα τα κομπιουτεράκια εκτελούν τις τέσσερις πράξεις. Για να κάνουμε μια πράξη μεταξύ δύο αριθμών, πληκτρολογούμε τον πρώτο, το σύμβολο της πράξης, τον δεύτερο και μετά το \key{$=$}. Τα σύμβολα είναι \kpl, \kmi, \kpo{} και \kdi. Το κόμμα των δεκαδικών είναι συνήθως τελεία στην οθόνη και στο πλήκτρο. Για παράδειγμα, για το $7{,}28:5{,}2$ πατάμε \key{7.28} \kdi\ \key{5.2} \kis{} και βλέπουμε $1{,}4$.

\tmimasmall{Η μνήμη}

Αν το κομπιουτεράκι έχει μνήμη, έχει τέσσερα πλήκτρα:
\begin{itemize}[leftmargin=1.2em,itemsep=0.3ex,topsep=2pt]
\item \key{MR} (ή \key{MRC}): εμφανίζει στην οθόνη τον αριθμό που βρίσκεται στη μνήμη.
\item \key{MC}: σβήνει το περιεχόμενο της μνήμης.
\item \key{M+}: προσθέτει στη μνήμη τον αριθμό της οθόνης.
\item \key{M$-$}: αφαιρεί από τη μνήμη τον αριθμό της οθόνης.
\end{itemize}
Η μνήμη είναι ό,τι χρειάζεται η παρένθεση του χαρτιού: φυλάμε ένα αποτέλεσμα, κάνουμε άλλο υπολογισμό και μετά τα συνδυάζουμε.

\tmimasmall{Παραστάσεις με πολλές πράξεις}

Δεν εμπιστευόμαστε τυφλά τη σειρά των πλήκτρων. Μερικά απλά κομπιουτεράκια εκτελούν τις πράξεις με τη σειρά που τις πληκτρολογούμε, χωρίς να υπολογίζουν πρώτα τους πολλαπλασιασμούς (τα επιστημονικά την σέβονται). Για παραστάσεις με παρενθέσεις ή μικτές πράξεις, είναι ασφαλές να κάνουμε:
\begin{enumerate}[leftmargin=1.6em,itemsep=0.3ex,topsep=2pt]
\item Μια πράξη τη φορά, με την προτεραιότητα του χαρτιού (παρενθέσεις, δυνάμεις, πολλαπλασιασμός/διαίρεση, μετά πρόσθεση/αφαίρεση).
\item Φύλαξη των αποτελεσμάτων που θα χρειαστούμε αργότερα με \key{M+} ή \key{M$-$}, ή απλή γραφή τους στο τετράδιο.
\end{enumerate}

\tmimasmall{Η οθόνη δείχνει προσέγγιση}

Η οθόνη έχει συγκεκριμένο πλήθος ψηφίων (συνήθως $8$ ή $10$). Όταν το αποτέλεσμα έχει περισσότερα ψηφία, το κομπιουτεράκι το κόβει ή το στρογγυλοποιεί. Γι' αυτό $1:3$ δείχνει $0{,}33333333$ και $2:3$ δείχνει $0{,}66666667$. Άρα $3\cdot0{,}33333333=0{,}99999999$ και όχι $1$. Δεν σημαίνει ότι το κομπιουτεράκι χάλασε, αλλά ότι δείχνει μια προσέγγιση.

\begin{prosoxi}
\begin{itemize}[leftmargin=1.2em,itemsep=0.3ex,topsep=2pt]
\item Το κομπιουτεράκι κάνει σωστά ό,τι του πεις, ακόμα και αν του πεις λάθος. Πριν από κάθε χρήση κάνω μια \textbf{εκτίμηση} με στρογγυλά νούμερα: τόσο περίπου πρέπει να βγει.
\item Δεν ξεχνάμε να σβήσουμε τη μνήμη (\key{MC}) πριν από νέο υπολογισμό, αλλιώς προσθέτουμε σε παλιό αριθμό.
\item Η τιμή που δείχνει η οθόνη μπορεί να είναι στρογγυλοποιημένη. Όταν χρειαζόμαστε ακρίβεια, δεν στρογγυλοποιούμε ενδιάμεσα αποτελέσματα, μόνο το τελικό.
\end{itemize}
\end{prosoxi}

\tmima{Λυμένα παραδείγματα}

\paradeigma{Το σύνολο του ταμείου με τη μνήμη}
\tixero Οι αποδείξεις είναι $12{,}35$, $7{,}8$, $23{,}45$, $9{,}99$ και $15{,}2$ ευρώ.\par
\tithelo Το σύνολο, και ο μέσος όρος ανά απόδειξη.\par
\possindeo Εκτίμηση: $12+8+23+10+15=68$. Πατάω \key{MC}, και για κάθε απόδειξη πληκτρολογώ τον αριθμό και \key{M+}: \key{12.35}\key{M+}\ \key{7.8}\key{M+}\ \key{23.45}\key{M+}\ \key{9.99}\key{M+}\ \key{15.2}\key{M+}. Μετά \key{MR}. Για τον μέσο όρο διαιρώ με $5$: $68{,}79:5=13{,}758$.\par
\apantisi Σύνολο $68{,}79$ €, κατά μέσο όρο $13{,}76$ € (στρογγυλοποιημένο στο λεπτό).\par
\elegxos Η εκτίμηση ($68$) είναι πολύ κοντά στο $68{,}79$, ενώ το $77{,}79$ της Πολυξένης είναι εκτός (η εκτίμηση είναι $68$).

\paradeigma{Σύνθετη παράσταση με μνήμη}
\tixero Η παράσταση $(2{,}4:0{,}6+1{,}5\cdot4)\cdot2{,}5-18{,}9:(0{,}7\cdot5+3{,}5:0{,}5)$.\par
\tithelo Να την υπολογίσω με κομπιουτεράκι που έχει μνήμη.\par
\possindeo Πρώτα η πρώτη παρένθεση: $2{,}4:0{,}6=4$, $1{,}5\cdot4=6$, $4+6=10$, $10\cdot2{,}5=25$. Πατάω \key{M+} (η μνήμη έχει $25$). Μετά η δεύτερη: $0{,}7\cdot5=3{,}5$, $3{,}5:0{,}5=7$, $3{,}5+7=10{,}5$ και $18{,}9:10{,}5=1{,}8$. Πατάω \key{M$-$} (η μνήμη έχει $25-1{,}8$). Τέλος \key{MR}.\par
\apantisi $23{,}2$.\par
\elegxos Εκτίμηση: $(4+6)\cdot2{,}5=25$ και $18{,}9:10{,}5\approx2$, άρα $\approx23$. Συμφωνεί.

\paradeigma{Ο έλεγχος με εκτίμηση}
\tixero Ο Ηλίας πληκτρολογεί $38{,}5\cdot4{,}2$ και η οθόνη δείχνει $16{,}17$.\par
\tithelo Να δω αν είναι σωστό.\par
\possindeo Εκτίμηση: $40\cdot4=160$. Το $16{,}17$ είναι δέκα φορές μικρότερο. Ο Ηλίας πάτησε $3{,}85$ αντί για $38{,}5$ (η υποδιαστολή σε λάθος θέση). Διορθώνω: $38{,}5\cdot4{,}2=161{,}7$.\par
\apantisi $161{,}7$.\par
\elegxos $161{,}7\approx160$. Και $3{,}85\cdot4{,}2=16{,}17$, το λάθος που είδαμε.

\paradeigma{Η οθόνη δείχνει προσέγγιση}
\tixero Το κομπιουτεράκι έχει οθόνη οκτώ ψηφίων.\par
\tithelo Τι δείχνει για το $2:3$ και πόσο βγαίνει το $3\cdot(2:3)$ αν χρησιμοποιήσω την οθόνη;\par
\possindeo $2:3=0{,}666\ldots$ και η οθόνη κρατάει οκτώ ψηφία, στρογγυλοποιώντας το τελευταίο: $0{,}66666667$. Τριπλασιάζω: $3\cdot0{,}66666667=2{,}00000001$.\par
\apantisi Η οθόνη δείχνει $0{,}66666667$, και το γινόμενο με το $3$ βγαίνει $2{,}00000001$ αντί για $2$.\par
\elegxos Η διαφορά $0{,}00000001$ είναι αμελητέα, αλλά δείχνει ότι δουλεύουμε με προσέγγιση.

\paradeigma{Πολλαπλασιασμοί που προστίθενται}
\tixero Στο παζάρι πουλήθηκαν $12$ κέικ των $2{,}8$ €, $8$ μπουκάλια των $1{,}5$ € και $20$ μπισκότα των $0{,}6$ €.\par
\tithelo Τα συνολικά έσοδα.\par
\possindeo Κάθε γινόμενο το προσθέτω στη μνήμη: \key{12}\kpo\key{2.8}\kis\key{M+}\ \key{8}\kpo\key{1.5}\kis\key{M+}\ \key{20}\kpo\key{0.6}\kis\key{M+}\ \key{MR}.\par
\apantisi $33{,}6+12+12=57{,}6$ €.\par
\elegxos Εκτίμηση: $12\cdot3=36$ και $8\cdot1{,}5=12$ και $20\cdot0{,}6=12$, άρα $\approx60$. Συμφωνεί.

\tmima{Ασκήσεις}

\zoni{Το κατάλαβα}{πλήκτρα, μνήμη, εκτίμηση}

\ex Υπολόγισε με κομπιουτεράκι: α) $318{,}76+47{,}092$ \quad β) $1\,000-87{,}654$ \quad γ) $24{,}6\cdot3{,}05$ \quad δ) $91{,}2:0{,}16$.

\ex Με εκτίμηση βρες ποιο αποτέλεσμα είναι σωστό χωρίς να κάνεις την πράξη: α) $48{,}7\cdot5{,}9$ είναι $28{,}733$ ή $287{,}33$ ή $2873{,}3$; β) $19{,}8:4{,}5$ είναι $0{,}44$ ή $4{,}4$ ή $44$;

\ex Γράψε με ποια πλήκτρα θα υπολογίσεις: α) $(4{,}5+2{,}7)\cdot1{,}25$ \quad β) $36:(0{,}9+2{,}1)$. Πόσο βγαίνουν;

\ex Ποια παράσταση υπολογίζεται και πόσο βγαίνει; α) \key{MC} \key{6.5}\kpo\key{4}\kis\key{M+} \key{3.2}\kpo\key{5}\kis\key{M$-$} \key{MR}. \quad β) \key{1.2}\kpl\key{0.8}\kis\kdi\key{0.5}\kis.

\ex Γράψε τα πλήκτρα που χρειάζονται για τον υπολογισμό της $2{,}5\cdot8+1{,}5\cdot6-4{,}25\cdot2$ με χρήση μνήμης. Πόσο είναι το αποτέλεσμα;

\ex Υπολόγισε με κομπιουτεράκι: α) $(3{,}6+4{,}4):0{,}8-2{,}5$ \quad β) $12{,}6:(4{,}2\cdot0{,}5)$ \quad γ) $0{,}25\cdot(10-3{,}2)+1{,}2$.

\ex Σωστό ή λάθος; α) Το κομπιουτεράκι δεν κάνει ποτέ λάθος. \quad β) Η οθόνη μετά το $1:3$ δείχνει την ακριβή τιμή του $\frac{1}{3}$. \quad γ) Το \key{M+} προσθέτει στη μνήμη τον αριθμό της οθόνης. \quad δ) Όλα τα κομπιουτεράκια εκτελούν πρώτα τους πολλαπλασιασμούς. \quad ε) Όταν έχω κομπιουτεράκι δεν χρειάζεται εκτίμηση.

\zoni{Το χρησιμοποιώ}{ταμείο, καύσιμα, μοιρασιά}

\ex \textbf{Το ταμείο.} Στο παζάρι έσοδα ήταν $57{,}6$ € και τα έξοδα (υλικά, χαρτικά) $18{,}35$ €. Πόσα καθαρά κέρδη έχει το ταμείο της εκδρομής; Έλεγξε με εκτίμηση.

\ex \textbf{Η απόδειξη.} Ο Στέλιος πληκτρολόγησε $6{,}75\cdot1{,}2$ και είδε $81$. Εξήγησε τι πάτησε λάθος και βρες το σωστό αποτέλεσμα.

\ex \textbf{Το ταξίδι.} Το αυτοκίνητο του θείου Λευτέρη καταναλώνει $7{,}8$ λίτρα ανά $100$ χλμ. Πόση βενζίνη θα χρειαστεί για μια διαδρομή $312$ χλμ.; Στρογγυλοποίησε στο δέκατο του λίτρου.

\ex \textbf{Η μοιρασιά.} Το ταμείο των $100$ € μοιράζεται εξίσου σε $7$ παιδιά. Τι δείχνει η οθόνη; Πόσα ευρώ παίρνει το καθένα στρογγυλοποιημένα στο λεπτό και πόσα θα έδινε συνολικά το ταμείο με αυτή τη στρογγυλοποίηση;

\ex \textbf{Το ρεύμα.} Μια λάμπα ισχύος $0{,}06$ kW ανάβει $4{,}5$ ώρες την ημέρα για $30$ ημέρες. Με τιμή $0{,}18$ € ανά kWh, πόσο κοστίζει στον μήνα; (Κόστος = ισχύς × ώρες × τιμή.)

\zoni{Το σκέφτομαι}{ο έλεγχος είναι δική μας δουλειά}

\ex Εξήγησε γιατί στο κομπιουτεράκι των $8$ ψηφίων το $3\cdot(1:3)$ δεν δίνει ακριβώς $1$.

\ex Με ένα απλό κομπιουτεράκι, αν πατήσεις \key{2}\kpl\key{3}\kpo\key{4}\kis{} βλέπεις συχνά $20$ και όχι $14$. Ποια σειρά εκτέλεσης ακολούθησε και πώς θα την πατήσεις για να βγει το σωστό;

\ex Χωρίς να κάνεις την πράξη, εξήγησε γιατί το $487\cdot0{,}98$ δεν μπορεί να είναι $4\,773{,}3$ και βρες τη σωστή τάξη μεγέθους.

\ex Πόσες φορές πρέπει να πατήσεις \kpl\ \key{0.1} \kis{} ξεκινώντας από το $0$, για να φτάσεις το $2{,}5$; Με ποια διαίρεση το βρίσκεις αμέσως;

\ex Κάνε τη διαίρεση $1:7$ στο κομπιουτεράκι σου. Γράψε όσα ψηφία φαίνονται και στρογγυλοποίησέ τον στα εκατοστά. Τι μοτίβο βλέπεις στα ψηφία;

\zoni{Για το σπίτι}{δοκίμασε το δικό σου κομπιουτεράκι}

\ex Πάρε τον λογαριασμό ρεύματος ή νερού του σπιτιού, και με το κομπιουτεράκι έλεγξε ένα από τα γινόμενα ή τα αθροίσματα που αναγράφονται. Βγαίνει το ίδιο;

\ex Βρες τις θερμοκρασίες της τελευταίας εβδομάδας και με το κομπιουτεράκι υπολόγισε τον μέσο όρο τους. Στρογγυλοποίησέ τον στο δέκατο.
